% ═══════════════════════════════════════════════════════════════
% METHODENHANDBUCH: AKTIVIERENDE UNTERRICHTSMETHODEN
% ═══════════════════════════════════════════════════════════════
% Sammlung erprobter Methoden für den Sprachunterricht
% Besonders geeignet für: Unterrichtsbesuche, Hospitationen
%
% Verwendung durch Claude:
%   - Dieses Dokument dient als Referenz für Materialerstellung
%   - Bei Anfrage "Erstelle Material für [METHODE]" nutze die
%     Spezifikation im Abschnitt "Material" der jeweiligen Methode
% ═══════════════════════════════════════════════════════════════

\documentclass[11pt, a4paper]{article}

% ═══════════════════════════════════════════════════════════════
% PAKETE
% ═══════════════════════════════════════════════════════════════

\usepackage[utf8]{inputenc}
\usepackage[T1]{fontenc}
\usepackage[ngerman]{babel}
\usepackage{helvet}
\renewcommand{\familydefault}{\sfdefault}

\usepackage[margin=2cm]{geometry}
\usepackage{enumitem}
\usepackage{tabularx}
\usepackage{booktabs}
\usepackage{longtable}
\usepackage{multirow}
\usepackage{array}

\usepackage{xcolor}
\usepackage{amssymb}
\usepackage{tcolorbox}
\tcbuselibrary{skins, breakable}

\usepackage{fancyhdr}
\usepackage{lastpage}
\usepackage{needspace}
\usepackage{hyperref}

% ═══════════════════════════════════════════════════════════════
% FARBEN
% ═══════════════════════════════════════════════════════════════

\definecolor{akzent}{HTML}{c41e3a}      % Spanisch-Karminrot
\definecolor{grau}{HTML}{888888}
\definecolor{hellgrau}{HTML}{f5f5f5}
\definecolor{erfolg}{HTML}{2a7a4b}
\definecolor{warnung}{HTML}{b35c00}
\definecolor{einstieg}{HTML}{2c5aa0}
\definecolor{wortschatz}{HTML}{7b1fa2}
\definecolor{grammatik}{HTML}{00695c}
\definecolor{sprechen}{HTML}{c41e3a}
\definecolor{sicherung}{HTML}{5d4037}

% ═══════════════════════════════════════════════════════════════
% KOPF- UND FUSSZEILE
% ═══════════════════════════════════════════════════════════════

\pagestyle{fancy}
\fancyhf{}
\renewcommand{\headrulewidth}{0.4pt}
\renewcommand{\footrulewidth}{0.4pt}

\lhead{Methodenhandbuch}
\chead{\textbf{Aktivierende Unterrichtsmethoden}}
\rhead{Sprachunterricht}

\lfoot{Stand: \today}
\cfoot{}
\rfoot{Seite \thepage\ von \pageref{LastPage}}

% ═══════════════════════════════════════════════════════════════
% BEFEHLE
% ═══════════════════════════════════════════════════════════════

% Phasen-Tags
\newcommand{\tagEinstieg}{\colorbox{einstieg}{\textcolor{white}{\footnotesize\textbf{EINSTIEG}}}}
\newcommand{\tagWortschatz}{\colorbox{wortschatz}{\textcolor{white}{\footnotesize\textbf{WORTSCHATZ}}}}
\newcommand{\tagGrammatik}{\colorbox{grammatik}{\textcolor{white}{\footnotesize\textbf{GRAMMATIK}}}}
\newcommand{\tagSprechen}{\colorbox{sprechen}{\textcolor{white}{\footnotesize\textbf{SPRECHEN}}}}
\newcommand{\tagSicherung}{\colorbox{sicherung}{\textcolor{white}{\footnotesize\textbf{SICHERUNG}}}}

% Zeit-Tag
\newcommand{\zeit}[1]{\textcolor{grau}{\footnotesize $\approx$ #1 Min.}}

% Beobachtungs-Eignung
\newcommand{\beobachtung}{\textcolor{erfolg}{\footnotesize $\bigstar$ \textbf{Beobachtung}}}

% Material-Marker
\newcommand{\matKeine}{\textcolor{erfolg}{$\checkmark$ kein Material}}
\newcommand{\matEinfach}{\textcolor{warnung}{$\circ$ einfaches Material}}
\newcommand{\matVorbereitung}{\textcolor{akzent}{$\bullet$ Vorbereitung nötig}}

% Methoden-Box
\newtcolorbox{methode}[2][]{
    enhanced,
    breakable,
    colback=white,
    colframe=akzent,
    colbacktitle=akzent,
    coltitle=white,
    boxrule=1pt,
    arc=0pt,
    fonttitle=\bfseries\large,
    title={#2},
    #1
}

% Material-Box
\newtcolorbox{materialbox}{
    colback=hellgrau,
    colframe=grau,
    boxrule=0.5pt,
    arc=0pt,
    left=5pt, right=5pt, top=5pt, bottom=5pt
}

% Tipp-Box
\newtcolorbox{tipp}{
    colback=white,
    colframe=erfolg,
    boxrule=0.5pt,
    arc=0pt,
    left=5pt, right=5pt, top=3pt, bottom=3pt,
    before upper={\textcolor{erfolg}{\textbf{Tipp:}} }
}

% ═══════════════════════════════════════════════════════════════
% DOKUMENT
% ═══════════════════════════════════════════════════════════════

\begin{document}

% ═══════════════════════════════════════════════════════════════
% TITELSEITE
% ═══════════════════════════════════════════════════════════════

\begin{center}
{\Huge\bfseries\textcolor{akzent}{Methodenhandbuch}}\\[0.5em]
{\LARGE Aktivierende Unterrichtsmethoden}\\[1em]
{\large für den Sprachunterricht}\\[2em]
{\normalsize Sammlung erprobter Methoden mit Kurzanleitung und Materialspezifikation}\\[1em]
{\footnotesize Stand: \today}
\end{center}

\vspace{2em}

% ═══════════════════════════════════════════════════════════════
% ÜBERSICHT
% ═══════════════════════════════════════════════════════════════

\section*{Übersicht nach Unterrichtsphase}

\renewcommand{\arraystretch}{1.3}
\begin{longtable}{|p{3.5cm}|c|c|c|p{4cm}|}
\hline
\textbf{Methode} & \textbf{Zeit} & \textbf{Material} & \textbf{UB} & \textbf{Phase} \\
\hline
\endfirsthead
\hline
\textbf{Methode} & \textbf{Zeit} & \textbf{Material} & \textbf{UB} & \textbf{Phase} \\
\hline
\endhead

\multicolumn{5}{l}{\textit{Fortsetzung auf nächster Seite...}} \\
\endfoot

\hline
\endlastfoot

% EINSTIEG
Kugellager & 10 & -- & $\bigstar$ & \tagEinstieg \\
\hline
Blitzlicht-Runde & 5 & -- & & \tagEinstieg \\
\hline
Bildimpuls & 5 & Bild & & \tagEinstieg \\
\hline

% WORTSCHATZ
Galeriegang & 15 & Karten & $\bigstar$ & \tagWortschatz \\
\hline
Korridor-Methode & 15 & Post-its & & \tagWortschatz \\
\hline
Vokabel-Staffel & 10 & Tafel & & \tagWortschatz \\
\hline
Flüsterpost & 5 & -- & & \tagWortschatz \\
\hline
Wortfeld-Mindmap & 10 & Tafel & & \tagWortschatz \\
\hline
Zuordnungs-Puzzle & 10 & Karten & & \tagWortschatz \\
\hline

% GRAMMATIK
Lebende Grammatik & 15 & Wortkarten & $\bigstar$ & \tagGrammatik \\
\hline
Fehler-Jagd & 10 & AB/Folie & & \tagGrammatik \\
\hline
Konjugations-Würfel & 10 & Würfel & & \tagGrammatik \\
\hline
Satzbausteine & 15 & Magnete & $\bigstar$ & \tagGrammatik \\
\hline

% SPRECHEN
Interview-Bingo & 15 & Bingo-Karte & $\bigstar$ & \tagSprechen \\
\hline
Fishbowl & 15 & -- & $\bigstar$ & \tagSprechen \\
\hline
Experten-Puzzle & 25 & Texte & $\bigstar$ & \tagSprechen \\
\hline
Rollenspiel & 15 & Rollenkarten & & \tagSprechen \\
\hline

% SICHERUNG
Daumen-Barometer & 2 & -- & & \tagSicherung \\
\hline
Exit Ticket & 3 & Post-its & $\bigstar$ & \tagSicherung \\
\hline
Drei-Wort-Zusammenfassung & 5 & -- & & \tagSicherung \\
\hline
Stille Post schriftlich & 10 & Papier & & \tagSicherung \\
\hline

\end{longtable}

\textbf{Legende:} UB = besonders geeignet für Unterrichtsbesuch ($\bigstar$)

\newpage

% ═══════════════════════════════════════════════════════════════
% EINSTIEG / WARM-UP
% ═══════════════════════════════════════════════════════════════

\section{Einstieg / Warm-up}

% --- KUGELLAGER ---
\begin{methode}{1. Kugellager / Speed Dating}

\tagEinstieg{} \zeit{8--12} \hfill \beobachtung

\subsection*{Didaktischer Hintergrund}
\textbf{Kooperatives Lernen} nach Norm \& Kathy Green. Maximale Sprechzeit für alle SuS gleichzeitig, niedrige Hemmschwelle durch wechselnde Partner.

\subsection*{Ablauf}
\begin{enumerate}[leftmargin=*]
    \item SuS bilden zwei konzentrische Kreise (Innenkreis schaut nach außen, Außenkreis nach innen)
    \item Paare stehen sich gegenüber
    \item L gibt Impuls/Frage (z.B. \textit{``¿Cómo te llamas?''})
    \item Innenkreis fragt, Außenkreis antwortet (30--45 Sek.)
    \item Signal $\rightarrow$ Außenkreis rotiert um eine Position
    \item Wiederholen mit neuen Partnern (5--8 Durchgänge)
\end{enumerate}

\subsection*{Material}
\begin{materialbox}
\begin{itemize}[label=$\square$, leftmargin=*]
    \item Kein Material zwingend erforderlich
    \item Optional: Fragekarten für schwächere SuS
    \item Optional: Timer/Gong für Signale
\end{itemize}

\textbf{Materialspezifikation für Erstellung:}
\begin{itemize}[leftmargin=*]
    \item \textbf{Fragekarten:} DIN A6, je eine Frage pro Karte, Spanisch + deutsche Hilfe auf Rückseite
    \item \textbf{Anzahl:} 8--10 verschiedene Fragen
    \item \textbf{Format:} Schneidevorlagen auf A4 (4 Karten pro Blatt)
\end{itemize}
\end{materialbox}

\subsection*{Varianten}
\begin{itemize}[leftmargin=*]
    \item \textbf{Stühle statt Kreise:} Zwei Stuhlreihen gegenüber (bei Platzmangel)
    \item \textbf{Themenwechsel:} Nach 4 Runden neue Fragestellung
    \item \textbf{Grammatik-Drill:} Konjugationsabfrage statt Fragen
\end{itemize}

\begin{tipp}
Bei ungerader SuS-Zahl: L übernimmt eine Position oder Dreiergruppe.
\end{tipp}

\end{methode}

\vspace{1em}

% --- BLITZLICHT ---
\begin{methode}{2. Blitzlicht-Runde}

\tagEinstieg{} \zeit{3--5} \hfill \matKeine

\subsection*{Didaktischer Hintergrund}
Schnelle Aktivierung des Vorwissens. Jede/r kommt zu Wort. Fokussierung auf das Thema.

\subsection*{Ablauf}
\begin{enumerate}[leftmargin=*]
    \item L nennt Thema/Oberbegriff (z.B. \textit{``la familia''})
    \item Jede/r SuS sagt EIN Wort zum Thema (reihum oder freiwillig)
    \item \textbf{Regel:} Keine Wiederholungen!
    \item Wird zunehmend schwieriger $\rightarrow$ Spannung
\end{enumerate}

\subsection*{Material}
\begin{materialbox}
Kein Material erforderlich.
\end{materialbox}

\subsection*{Varianten}
\begin{itemize}[leftmargin=*]
    \item \textbf{Mit Ball:} Ball werfen statt Reihenfolge
    \item \textbf{Satz statt Wort:} \textit{``Tengo un hermano.''}
    \item \textbf{Ausscheiden:} Wer kein Wort findet, setzt sich (spielerisch)
\end{itemize}

\end{methode}

\vspace{1em}

% --- BILDIMPULS ---
\begin{methode}{3. Bildimpuls mit Vermutung}

\tagEinstieg{} \zeit{5--8} \hfill \matEinfach

\subsection*{Didaktischer Hintergrund}
Visuelle Aktivierung, Neugier wecken, freies Sprechen anregen. Nutzt natürliche Neugierde.

\subsection*{Ablauf}
\begin{enumerate}[leftmargin=*]
    \item Zeige Bild (zunächst verdeckt/abgeschnitten/verpixelt)
    \item SuS vermuten: \textit{``¿Qué es?'' / ``¿Qué pasa?''}
    \item Schrittweise mehr zeigen
    \item Auflösung $\rightarrow$ Überleitung zum Thema
\end{enumerate}

\subsection*{Material}
\begin{materialbox}
\begin{itemize}[label=$\square$, leftmargin=*]
    \item 1 Bild (Beamer/Folie/Ausdruck)
    \item Abdeckmaterial oder digitale Zoom-Funktion
\end{itemize}

\textbf{Bildauswahl-Kriterien:}
\begin{itemize}[leftmargin=*]
    \item Thematisch passend zur Stunde
    \item Mehrdeutig genug für Vermutungen
    \item Kulturell relevant (z.B. hispanische Länder)
\end{itemize}
\end{materialbox}

\end{methode}

\newpage

% ═══════════════════════════════════════════════════════════════
% WORTSCHATZARBEIT
% ═══════════════════════════════════════════════════════════════

\section{Wortschatzarbeit}

% --- GALERIEGANG ---
\begin{methode}{4. Galeriegang / Gallery Walk}

\tagWortschatz{} \zeit{12--15} \hfill \beobachtung

\subsection*{Didaktischer Hintergrund}
Bewegtes Lernen, selbstgesteuertes Tempo, visuelle Verankerung. Basiert auf dem Prinzip der \textbf{Lernstationen} kombiniert mit \textbf{Bewegungslernen}.

\subsection*{Ablauf}
\begin{enumerate}[leftmargin=*]
    \item Vokabelkarten (Wort + Bild, KEINE Übersetzung) im Raum verteilen
    \item SuS gehen umher, betrachten Karten
    \item Notieren: Welche Wörter erkenne ich? (in Tabelle)
    \item Nach 10 Min: Partnervergleich
    \item Plenum: Unbekannte Wörter klären
\end{enumerate}

\subsection*{Material}
\begin{materialbox}
\begin{itemize}[label=$\square$, leftmargin=*]
    \item 15--20 Vokabelkarten (DIN A5 oder A4)
    \item Klebeband/Magnete zum Befestigen
    \item Notizzettel für SuS
\end{itemize}

\textbf{Materialspezifikation für Erstellung:}
\begin{itemize}[leftmargin=*]
    \item \textbf{Vokabelkarten:} DIN A5, Wort groß + passendes Bild/Icon
    \item \textbf{Layout:} Wort oben (Spanisch), Bild darunter
    \item \textbf{KEINE Übersetzung auf Vorderseite}
    \item \textbf{Optional:} QR-Code auf Rückseite $\rightarrow$ Aussprache-Audio
    \item \textbf{Notizblatt:} Tabelle mit 2 Spalten (Spanisch | Vermutung)
\end{itemize}
\end{materialbox}

\begin{tipp}
Karten auf Augenhöhe anbringen. Mix aus bekannten und neuen Wörtern (ca. 60/40).
\end{tipp}

\end{methode}

\vspace{1em}

% --- KORRIDOR-METHODE ---
\begin{methode}{5. Korridor-Methode}

\tagWortschatz{} \zeit{10--15} \hfill \matVorbereitung

\subsection*{Didaktischer Hintergrund}
Bewegung + Eigenverantwortung. SuS bestimmen Lerntempo selbst. Multisensorisch: Sehen, Gehen, Anfassen.

\subsection*{Ablauf}
\begin{enumerate}[leftmargin=*]
    \item Post-its mit Vokabeln entlang des Korridors/Raumes auslegen
    \item Vorderseite: Spanisches Wort
    \item Rückseite: Deutsche Übersetzung
    \item SuS gehen einzeln, nehmen Post-it auf, lernen, legen zurück
    \item Weitergehen zum nächsten
    \item Nach 10 Min: Schnelltest (L zeigt Wort, SuS übersetzen)
\end{enumerate}

\subsection*{Material}
\begin{materialbox}
\begin{itemize}[label=$\square$, leftmargin=*]
    \item Post-its (groß, ca. 7$\times$7 cm)
    \item Oder: Laminierte Karten (wiederverwendbar)
    \item Edding für gut lesbare Schrift
\end{itemize}

\textbf{Materialspezifikation für Erstellung:}
\begin{itemize}[leftmargin=*]
    \item \textbf{Druckvorlage:} 2-seitig, Vorderseite ES / Rückseite DE
    \item \textbf{Format:} 6 Karten pro A4-Blatt
    \item \textbf{Schriftgröße:} mind. 24pt
    \item \textbf{Anzahl:} 15--25 Vokabeln pro Durchgang
\end{itemize}
\end{materialbox}

\subsection*{Varianten}
\begin{itemize}[leftmargin=*]
    \item \textbf{Partnerversion:} Zu zweit gehen, gegenseitig abfragen
    \item \textbf{Wettbewerb:} Wer kann am Ende die meisten?
\end{itemize}

\end{methode}

\vspace{1em}

% --- VOKABEL-STAFFEL ---
\begin{methode}{6. Vokabel-Staffel}

\tagWortschatz{} \zeit{8--10} \hfill \matKeine

\subsection*{Didaktischer Hintergrund}
Gamification, Teamwettbewerb, Bewegung. Aktiviert Wortschatz unter Zeitdruck.

\subsection*{Ablauf}
\begin{enumerate}[leftmargin=*]
    \item Klasse in 3--4 Teams teilen
    \item Jedes Team steht in einer Reihe vor der Tafel
    \item L gibt Kategorie vor (z.B. \textit{``colores''})
    \item Auf Signal: Erste/r jedes Teams läuft zur Tafel, schreibt ein Wort
    \item Zurück, abklatschen, nächste/r läuft
    \item \textbf{Regel:} Keine Wiederholungen!
    \item Erstes Team mit 10 korrekten Wörtern gewinnt
\end{enumerate}

\subsection*{Material}
\begin{materialbox}
\begin{itemize}[label=$\square$, leftmargin=*]
    \item Tafel/Whiteboard (in Sektoren unterteilt)
    \item Kreide/Marker (pro Team eine Farbe)
\end{itemize}
\end{materialbox}

\begin{tipp}
Orthografie-Fehler = Wort zählt nicht! Motiviert zur Sorgfalt.
\end{tipp}

\end{methode}

\vspace{1em}

% --- ZUORDNUNGS-PUZZLE ---
\begin{methode}{7. Zuordnungs-Puzzle / Satz-Hälften}

\tagWortschatz{} \zeit{8--10} \hfill \matVorbereitung

\subsection*{Didaktischer Hintergrund}
Bewegung + Sozialform-Wechsel. Kooperatives Finden des Partners. Sinnerfassendes Lesen.

\subsection*{Ablauf}
\begin{enumerate}[leftmargin=*]
    \item Sätze in zwei Hälften schneiden
    \item Karten mischen, an SuS verteilen (jede/r eine Hälfte)
    \item SuS suchen ihre ``passende Hälfte'' im Raum
    \item Gefundene Paare lesen Satz vor
    \item Klasse validiert: Korrekt?
\end{enumerate}

\subsection*{Material}
\begin{materialbox}
\begin{itemize}[label=$\square$, leftmargin=*]
    \item Satzkarten (halbiert)
    \item Anzahl: Klassengröße / 2
\end{itemize}

\textbf{Materialspezifikation für Erstellung:}
\begin{itemize}[leftmargin=*]
    \item \textbf{Sätze:} Thematisch passend, grammatisch korrekt
    \item \textbf{Schnittposition:} Logische Trennstelle (nach Subjekt/Verb)
    \item \textbf{Format:} Streifen (ca. 20$\times$5 cm)
    \item \textbf{Differenzierung:} Farbliche Markierung für Schwierigkeitsgrade
\end{itemize}
\end{materialbox}

\end{methode}

\newpage

% ═══════════════════════════════════════════════════════════════
% GRAMMATIK
% ═══════════════════════════════════════════════════════════════

\section{Grammatik-Einführung und Übung}

% --- LEBENDE GRAMMATIK ---
\begin{methode}{8. Lebende Grammatik / Human Sentences}

\tagGrammatik{} \zeit{12--15} \hfill \beobachtung

\subsection*{Didaktischer Hintergrund}
\textbf{Total Physical Response (TPR)} nach James Asher. Körperliche Bewegung verankert Strukturen. Visualisiert Satzstellung.

\subsection*{Ablauf}
\begin{enumerate}[leftmargin=*]
    \item Wortkarten vorbereiten (je ein Wort pro Karte)
    \item SuS erhalten je eine Karte
    \item Aufgabe: ``Bildet den Satz: Yo como una manzana.''
    \item SuS mit passenden Karten stellen sich in Reihenfolge auf
    \item Klasse prüft: Korrekte Reihenfolge?
    \item Variante: Negation einfügen $\rightarrow$ ``no'' findet seinen Platz
\end{enumerate}

\subsection*{Material}
\begin{materialbox}
\begin{itemize}[label=$\square$, leftmargin=*]
    \item Wortkarten (DIN A4, gut lesbar)
    \item Ca. 15--20 Wörter (Subjekte, Verben, Objekte, Artikel)
\end{itemize}

\textbf{Materialspezifikation für Erstellung:}
\begin{itemize}[leftmargin=*]
    \item \textbf{Format:} DIN A4 quer, ein Wort pro Karte
    \item \textbf{Schriftgröße:} mind. 72pt, fett
    \item \textbf{Farbcodierung optional:}
    \begin{itemize}
        \item Subjekte: blau
        \item Verben: rot
        \item Objekte: grün
        \item Artikel/Präpositionen: grau
    \end{itemize}
    \item \textbf{Wörter für Spanisch-Grundlagen:}
    \begin{itemize}
        \item Subjekte: yo, tú, él, ella, nosotros, ellos
        \item Verben: como, bebo, hablo, tengo, soy, estoy
        \item Objekte: manzana, agua, libro, hermano, casa
        \item Artikel: un, una, el, la, los, las
        \item Negation: no
        \item Adjektive: grande, pequeño, nuevo
    \end{itemize}
\end{itemize}
\end{materialbox}

\begin{tipp}
Magnete auf Rückseite $\rightarrow$ Karten auch an Tafel verwendbar.
\end{tipp}

\end{methode}

\vspace{1em}

% --- FEHLER-JAGD ---
\begin{methode}{9. Fehler-Jagd / Error Hunt}

\tagGrammatik{} \zeit{8--10} \hfill \matEinfach

\subsection*{Didaktischer Hintergrund}
Aktiviert kritisches Lesen. SuS wenden Regelwissen an. Spielerischer Wettbewerb.

\subsection*{Ablauf}
\begin{enumerate}[leftmargin=*]
    \item Text mit 5--7 absichtlichen Fehlern zeigen (Beamer/Folie)
    \item SuS arbeiten in Paaren
    \item Aufgabe: ``Findet alle Fehler und korrigiert sie!''
    \item Erstes Paar mit allen korrekten Lösungen gewinnt
    \item Gemeinsame Besprechung: Warum ist es falsch?
\end{enumerate}

\subsection*{Material}
\begin{materialbox}
\begin{itemize}[label=$\square$, leftmargin=*]
    \item Fehlertext (Folie/Beamer/AB)
    \item Lösungsfolie
\end{itemize}

\textbf{Materialspezifikation für Erstellung:}
\begin{itemize}[leftmargin=*]
    \item \textbf{Textlänge:} 5--8 Sätze
    \item \textbf{Fehlertypen:} Mix aus
    \begin{itemize}
        \item Genus-Fehler (*el casa)
        \item Konjugationsfehler (*yo comes)
        \item Akzent-Fehler (*esta statt está)
        \item Wortstellungsfehler
    \end{itemize}
    \item \textbf{Schwierigkeit:} Fehler erkennbar, aber nicht trivial
    \item \textbf{Lösungsblatt:} Fehler markiert + Korrektur + Erklärung
\end{itemize}
\end{materialbox}

\end{methode}

\vspace{1em}

% --- KONJUGATIONS-WÜRFEL ---
\begin{methode}{10. Konjugations-Würfel}

\tagGrammatik{} \zeit{8--10} \hfill \matVorbereitung

\subsection*{Didaktischer Hintergrund}
Drill mit Gamification. Zufallselement erhöht Aufmerksamkeit. Schnelles Tempo = Automatisierung.

\subsection*{Ablauf}
\begin{enumerate}[leftmargin=*]
    \item Würfel mit Personalpronomen (yo, tú, él, nosotros, vosotros, ellos)
    \item L nennt Infinitiv (z.B. ``hablar'')
    \item SuS würfelt $\rightarrow$ zeigt z.B. ``nosotros''
    \item SuS muss konjugieren: ``hablamos''
    \item Korrekt = Punkt / Falsch = nächste/r
\end{enumerate}

\subsection*{Material}
\begin{materialbox}
\begin{itemize}[label=$\square$, leftmargin=*]
    \item 1--2 Würfel mit Pronomen
    \item Alternative: Normaler Würfel + Zuordnungsliste (1=yo, 2=tú, ...)
\end{itemize}

\textbf{Materialspezifikation für Erstellung:}
\begin{itemize}[leftmargin=*]
    \item \textbf{Würfel-Vorlage:} Faltvorlage für Papier-Würfel
    \item \textbf{Seiten:} yo, tú, él/ella, nosotros, vosotros, ellos/ellas
    \item \textbf{Variante 2:} Verb-Würfel + Pronomen-Würfel = 2 würfeln
\end{itemize}
\end{materialbox}

\end{methode}

\vspace{1em}

% --- SATZBAUSTEINE ---
\begin{methode}{11. Satzbausteine / Magnetkarten}

\tagGrammatik{} \zeit{12--15} \hfill \beobachtung

\subsection*{Didaktischer Hintergrund}
Visualisierung von Satzstruktur. Haptisches Manipulieren. Fehlerkorrektur durch Umstellen sichtbar.

\subsection*{Ablauf}
\begin{enumerate}[leftmargin=*]
    \item Laminierte Wortkarten mit Magneten an Tafel
    \item L oder SuS bildet (falschen) Satz
    \item Klasse korrigiert durch Umstellen
    \item Diskussion: Warum diese Reihenfolge?
\end{enumerate}

\subsection*{Material}
\begin{materialbox}
\begin{itemize}[label=$\square$, leftmargin=*]
    \item Laminierte Wortkarten (ca. 10$\times$5 cm)
    \item Magnetstreifen auf Rückseite
    \item Magnettafel
\end{itemize}

\textbf{Materialspezifikation für Erstellung:}
\begin{itemize}[leftmargin=*]
    \item \textbf{Format:} Streifen, ca. 15$\times$4 cm
    \item \textbf{Laminiert} für Wiederverwendung
    \item \textbf{Wortarten farbcodiert} (siehe ``Lebende Grammatik'')
    \item \textbf{Set pro Thema:} ca. 25--30 Karten
\end{itemize}
\end{materialbox}

\end{methode}

\newpage

% ═══════════════════════════════════════════════════════════════
% SPRECHEN / KOMMUNIKATION
% ═══════════════════════════════════════════════════════════════

\section{Sprechen / Kommunikation}

% --- INTERVIEW-BINGO ---
\begin{methode}{12. Interview-Bingo}

\tagSprechen{} \zeit{12--15} \hfill \beobachtung

\subsection*{Didaktischer Hintergrund}
Authentische Kommunikationsanlässe. Alle SuS gleichzeitig aktiv. Bewegung im Raum. Ziel-orientiertes Fragen.

\subsection*{Ablauf}
\begin{enumerate}[leftmargin=*]
    \item Jede/r SuS erhält Bingo-Karte mit Aussagen
    \item Beispiel: ``Tiene un hermano'' / ``Le gusta el fútbol''
    \item SuS gehen umher, stellen Fragen auf Spanisch
    \item Passt die Antwort $\rightarrow$ Name ins Feld
    \item Erste/r mit vollständiger Reihe ruft ``¡Bingo!''
\end{enumerate}

\subsection*{Material}
\begin{materialbox}
\begin{itemize}[label=$\square$, leftmargin=*]
    \item Bingo-Karten (DIN A5 oder A4)
    \item 9--16 Felder (3$\times$3 oder 4$\times$4)
\end{itemize}

\textbf{Materialspezifikation für Erstellung:}
\begin{itemize}[leftmargin=*]
    \item \textbf{Format:} A5, Raster 3$\times$3 oder 4$\times$4
    \item \textbf{Aussagen:} 3. Person Singular
    \item \textbf{Thematisch:} Passend zur aktuellen Einheit
    \item \textbf{Variation:} 2--3 verschiedene Karten (andere Anordnung)
    \item \textbf{Beispiel-Aussagen (Familia):}
    \begin{itemize}
        \item Tiene un hermano / Tiene una hermana
        \item Tiene un perro / Tiene un gato
        \item Vive con sus abuelos
        \item Es hijo único / Es hija única
    \end{itemize}
    \item \textbf{Fragehilfe auf Rückseite:} ``¿Tienes...?'' / ``¿Vives...?''
\end{itemize}
\end{materialbox}

\begin{tipp}
Differenzierung: Schwächere SuS bekommen Karte mit Fragehilfen.
\end{tipp}

\end{methode}

\vspace{1em}

% --- FISHBOWL ---
\begin{methode}{13. Fishbowl-Diskussion}

\tagSprechen{} \zeit{15--20} \hfill \beobachtung

\subsection*{Didaktischer Hintergrund}
Fokussierte Aufmerksamkeit. Modellieren von gutem Sprechen. Beobachtungsauftrag für Außenkreis.

\subsection*{Ablauf}
\begin{enumerate}[leftmargin=*]
    \item Innenkreis: 4--5 SuS (Diskutierende)
    \item Außenkreis: Rest der Klasse (Beobachter)
    \item Diskussionsimpuls: z.B. ``¿Qué es mejor: vivir en la ciudad o en el campo?''
    \item Innenkreis diskutiert 5--7 Minuten
    \item Außenkreis beobachtet mit Beobachtungsbogen
    \item Wechsel: Neue SuS in den Innenkreis
    \item Abschluss: Feedback zur Sprachverwendung
\end{enumerate}

\subsection*{Material}
\begin{materialbox}
\begin{itemize}[label=$\square$, leftmargin=*]
    \item Diskussionsimpuls / Fragestellung
    \item Beobachtungsbogen für Außenkreis
    \item Optional: Redemittel-Karten
\end{itemize}

\textbf{Materialspezifikation für Erstellung:}
\begin{itemize}[leftmargin=*]
    \item \textbf{Beobachtungsbogen:}
    \begin{itemize}
        \item Name des Beobachteten
        \item Kriterien: Verwendete Strukturen, Fehler, Flüssigkeit
        \item Notizfeld für gute Formulierungen
    \end{itemize}
    \item \textbf{Redemittel-Karte:}
    \begin{itemize}
        \item Zustimmung: Estoy de acuerdo, Tienes razón
        \item Widerspruch: No estoy de acuerdo, Pienso que...
        \item Nachfragen: ¿Puedes repetir?, ¿Qué quieres decir?
    \end{itemize}
\end{itemize}
\end{materialbox}

\end{methode}

\vspace{1em}

% --- EXPERTEN-PUZZLE ---
\begin{methode}{14. Experten-Puzzle / Jigsaw}

\tagSprechen{} \zeit{20--25} \hfill \beobachtung

\subsection*{Didaktischer Hintergrund}
\textbf{Kooperatives Lernen} nach Aronson. Positive Interdependenz: Jede/r ist Expert/in für einen Teil. Lehren durch Erklären.

\subsection*{Ablauf}
\begin{enumerate}[leftmargin=*]
    \item \textbf{Phase 1 -- Stammgruppen:} 4er-Gruppen, jede/r bekommt Text A/B/C/D
    \item \textbf{Phase 2 -- Expertengruppen:} Alle A's treffen sich, alle B's, etc.
    \item Expertengruppen erarbeiten ihren Text gemeinsam
    \item \textbf{Phase 3 -- Rückkehr:} Zurück in Stammgruppen
    \item Jede/r erklärt seinen Teil den anderen
    \item \textbf{Phase 4 -- Sicherung:} Gemeinsame Aufgabe / Quiz
\end{enumerate}

\subsection*{Material}
\begin{materialbox}
\begin{itemize}[label=$\square$, leftmargin=*]
    \item 4 verschiedene Texte/Infoblätter (A, B, C, D)
    \item Leitfragen für Expertengruppen
    \item Abschluss-Quiz oder Arbeitsblatt
\end{itemize}

\textbf{Materialspezifikation für Erstellung:}
\begin{itemize}[leftmargin=*]
    \item \textbf{Texte:} Gleiche Länge, gleicher Schwierigkeitsgrad
    \item \textbf{Kennzeichnung:} Große Buchstaben A/B/C/D oben
    \item \textbf{Leitfragen:} 3--4 Fragen pro Text
    \item \textbf{Format:} A4, Textlänge ca. 100--150 Wörter
    \item \textbf{Abschluss-AB:} Fragen zu allen 4 Teilen (nur lösbar, wenn alle beitragen)
\end{itemize}
\end{materialbox}

\begin{tipp}
Timer für jede Phase. Rollenverteilung in Expertengruppe: Vorleser, Schreiber, Zeitwächter.
\end{tipp}

\end{methode}

\vspace{1em}

% --- ROLLENSPIEL ---
\begin{methode}{15. Rollenspiel mit Rollenkarten}

\tagSprechen{} \zeit{12--15} \hfill \matVorbereitung

\subsection*{Didaktischer Hintergrund}
Authentische Kommunikationssituation. Identifikation mit Rolle senkt Hemmschwelle. Zielorientiertes Sprechen.

\subsection*{Ablauf}
\begin{enumerate}[leftmargin=*]
    \item Paare oder Kleingruppen bilden
    \item Jede/r erhält Rollenkarte mit:
    \begin{itemize}
        \item Charakter-Info
        \item Geheimes Ziel
    \end{itemize}
    \item Vorbereitungszeit: 2--3 Minuten
    \item Durchführung des Rollenspiels
    \item Optional: Präsentation vor der Klasse
\end{enumerate}

\subsection*{Material}
\begin{materialbox}
\begin{itemize}[label=$\square$, leftmargin=*]
    \item Rollenkarten (2 pro Paar)
    \item Situationsbeschreibung
\end{itemize}

\textbf{Materialspezifikation für Erstellung:}
\begin{itemize}[leftmargin=*]
    \item \textbf{Format:} DIN A6 oder halbes A5
    \item \textbf{Inhalt pro Karte:}
    \begin{itemize}
        \item Situation (kurz)
        \item Rolle (Name, Alter, Kontext)
        \item Geheimes Ziel (das der Partner nicht weiß)
    \end{itemize}
    \item \textbf{Beispiel Restaurant:}
    \begin{itemize}
        \item Karte A: Du bist Kellner. Empfiehl das Tagesgericht.
        \item Karte B: Du bist Gast. Du hast nur 10€. Frag nach günstigen Optionen.
    \end{itemize}
    \item \textbf{Hilfen:} Relevante Redemittel auf Rückseite
\end{itemize}
\end{materialbox}

\end{methode}

\newpage

% ═══════════════════════════════════════════════════════════════
% SICHERUNG / ABSCHLUSS
% ═══════════════════════════════════════════════════════════════

\section{Sicherung / Abschluss}

% --- DAUMEN-BAROMETER ---
\begin{methode}{16. Daumen-Barometer}

\tagSicherung{} \zeit{2--3} \hfill \matKeine

\subsection*{Didaktischer Hintergrund}
Schnelle Selbsteinschätzung. Visuelle Rückmeldung für L. Keine Bloßstellung einzelner SuS.

\subsection*{Ablauf}
\begin{enumerate}[leftmargin=*]
    \item L formuliert Kann-Aussage: ``Ich kann jetzt die Uhrzeit auf Spanisch sagen.''
    \item SuS zeigen gleichzeitig:
    \begin{itemize}
        \item Daumen hoch = Ja, sicher
        \item Daumen waagerecht = Teilweise / unsicher
        \item Daumen runter = Nein, noch nicht
    \end{itemize}
    \item L notiert sich Tendenz für Folgestunde
\end{enumerate}

\subsection*{Material}
\begin{materialbox}
Kein Material erforderlich.

\textbf{Alternative:} Ampelkarten (rot/gelb/grün) für weniger auffällige Rückmeldung.
\end{materialbox}

\end{methode}

\vspace{1em}

% --- EXIT TICKET ---
\begin{methode}{17. Exit Ticket / Eintrittskarte}

\tagSicherung{} \zeit{3--5} \hfill \beobachtung

\subsection*{Didaktischer Hintergrund}
Schriftliche Reflexion. Individuelles Feedback für L. Sammlung zeigt Lernstand der Klasse.

\subsection*{Ablauf}
\begin{enumerate}[leftmargin=*]
    \item 3 Minuten vor Stundenende
    \item Jede/r SuS schreibt auf Post-it:
    \begin{itemize}
        \item Variante A: ``Hoy he aprendido...''
        \item Variante B: Eine Frage, die noch offen ist
        \item Variante C: Ein neues Wort + Beispielsatz
    \end{itemize}
    \item Beim Hinausgehen: Post-it an Tür/Tafel kleben
    \item L sammelt, liest, reagiert in Folgestunde
\end{enumerate}

\subsection*{Material}
\begin{materialbox}
\begin{itemize}[label=$\square$, leftmargin=*]
    \item Post-its (Standardgröße)
    \item Sammelort (Tür, Plakat, Box)
\end{itemize}
\end{materialbox}

\begin{tipp}
Folgestunde mit Antworten auf Exit-Ticket-Fragen beginnen $\rightarrow$ zeigt Wertschätzung.
\end{tipp}

\end{methode}

\vspace{1em}

% --- DREI-WORT ---
\begin{methode}{18. Drei-Wort-Zusammenfassung}

\tagSicherung{} \zeit{3--5} \hfill \matKeine

\subsection*{Didaktischer Hintergrund}
Zwingt zur Synthese. Fokussiert auf Wesentliches. Schnell und aktivierend.

\subsection*{Ablauf}
\begin{enumerate}[leftmargin=*]
    \item L: ``Fasst die Stunde in genau DREI Wörtern zusammen.''
    \item 30 Sekunden Denkzeit
    \item Reihum oder per Zufall: Jede/r sagt seine drei Wörter
    \item L sammelt an Tafel $\rightarrow$ Wortwolke entsteht
\end{enumerate}

\subsection*{Material}
\begin{materialbox}
Kein Material erforderlich.

\textbf{Variante:} Digital mit Mentimeter o.ä. für Wortwolken-Visualisierung.
\end{materialbox}

\end{methode}

\vspace{1em}

% --- STILLE POST SCHRIFTLICH ---
\begin{methode}{19. Stille Post schriftlich}

\tagSicherung{} \zeit{8--10} \hfill \matEinfach

\subsection*{Didaktischer Hintergrund}
Spielerisch, kreativ, entspannend. Ergebnisse oft humorvoll $\rightarrow$ positiver Stundenabschluss.

\subsection*{Ablauf}
\begin{enumerate}[leftmargin=*]
    \item Jede/r SuS hat ein Blatt
    \item Erste/r schreibt einen Satz zum Thema
    \item Blatt falten, sodass nur das letzte Wort sichtbar
    \item Weitergeben an Nachbar/in
    \item Nächste/r schreibt Satz, der mit diesem Wort beginnt
    \item Wieder falten, weitergeben (5--6 Runden)
    \item Auffalten, vorlesen $\rightarrow$ oft absurde Geschichten
\end{enumerate}

\subsection*{Material}
\begin{materialbox}
\begin{itemize}[label=$\square$, leftmargin=*]
    \item Papierstreifen (längs halbes A4)
    \item Oder: Normales Papier, mehrfach gefaltet
\end{itemize}
\end{materialbox}

\end{methode}

\newpage

% ═══════════════════════════════════════════════════════════════
% ANHANG: BEOBACHTUNGS-CHECKLISTE
% ═══════════════════════════════════════════════════════════════

\section*{Anhang: Methodenwahl für Unterrichtsbesuche}

\subsection*{Checkliste: Ist die Methode beobachtungstauglich?}

\begin{itemize}[label=$\square$, leftmargin=*]
    \item[$\square$] \textbf{Sichtbare Schüleraktivierung:} Alle SuS sind erkennbar aktiv (nicht nur 2--3)
    \item[$\square$] \textbf{Klare Struktur:} Ablauf ist für Beobachter nachvollziehbar
    \item[$\square$] \textbf{Zeitlich planbar:} Methode lässt sich in definierten Zeitrahmen durchführen
    \item[$\square$] \textbf{Differenzierung erkennbar:} Verschiedene Niveaus sichtbar
    \item[$\square$] \textbf{Kompetenzbezug:} Fördert klar benennbare Kompetenz (Sprechen, Schreiben, ...)
    \item[$\square$] \textbf{Ergebnissicherung:} Lernprodukt am Ende sichtbar/dokumentierbar
\end{itemize}

\subsection*{Top 5 für den Unterrichtsbesuch}

\begin{tabularx}{\textwidth}{|c|l|X|l|}
\hline
\textbf{\#} & \textbf{Methode} & \textbf{Zeigt} & \textbf{Phase} \\
\hline
1 & Kugellager & Max. Sprechzeit, Kooperatives Lernen & Einstieg \\
\hline
2 & Interview-Bingo & Authentische Kommunikation, Bewegung & Übung \\
\hline
3 & Lebende Grammatik & Handlungsorientierung, Visualisierung & Erarbeitung \\
\hline
4 & Experten-Puzzle & Kooperatives Lernen, Selbstständigkeit & Erarbeitung \\
\hline
5 & Exit Ticket & Lernzielkontrolle, Reflexion & Sicherung \\
\hline
\end{tabularx}

\vspace{2em}

\subsection*{Material-Erstellungsaufträge an Claude}

Nutze folgende Formulierungen, um passendes Material zu erhalten:

\begin{tcolorbox}[colback=hellgrau, colframe=grau, boxrule=0.5pt, arc=0pt]
\textbf{Beispiel-Prompts:}
\begin{itemize}[leftmargin=*]
    \item ``Erstelle Fragekarten für Kugellager zum Thema \textit{[THEMA]}, Klasse \textit{[X]}.''
    \item ``Erstelle eine Interview-Bingo-Karte zum Thema \textit{[THEMA]} mit 9 Feldern.''
    \item ``Erstelle Wortkarten für Lebende Grammatik: Konjugation \textit{[VERB]}.''
    \item ``Erstelle einen Fehlertext zum Thema \textit{[GRAMMATIK]} mit 6 Fehlern.''
    \item ``Erstelle Rollenkarten für ein Rollenspiel im Restaurant.''
    \item ``Erstelle 4 Jigsaw-Texte zum Thema \textit{[LANDESKUNDE]}.''
\end{itemize}
\end{tcolorbox}

\vspace{2em}

\begin{center}
\textit{--- Ende des Methodenhandbuchs ---}
\end{center}

\end{document}
